\chapter{Foot Pain}

\section*{Definition}
Pain or discomfort localized anywhere in the foot, from the hindfoot (heel, ankle) through the midfoot (arch) to the forefoot (toes, metatarsals), with causes ranging from mechanical to systemic.

\section*{Pathophysiology}
Foot pain etiology varies by location: heel pain most commonly from plantar fasciitis; midfoot pain from degenerative arthritis or stress fracture; metatarsalgia (forefoot pain) from weight-bearing overload, Morton neuroma, or Freiberg infarction; great toe pain from hallux valgus (bunion), hallux rigidus, or gout. The foot's 26 bones, 33 joints, and complex ligamentous/tendinous architecture create numerous pain generators.

\section*{Etiology}
\begin{itemize}
  \item Plantar fasciitis
  \item Metatarsalgia
  \item Morton neuroma (interdigital neuroma)
  \item Hallux valgus (bunion)
  \item Hallux rigidus (great toe arthritis)
  \item Hammer toe or claw toe deformity
  \item Stress fracture (metatarsal, navicular)
  \item Tarsal tunnel syndrome
  \item Gout (usually first metatarsophalangeal joint)
  \item Flatfoot (pes planus) or cavus foot deformity
\end{itemize}

\section*{Indications for Evaluation}
Seek care if:
\begin{itemize}
  \item Severe or worsening symptoms
  \item Foot pain with inability to bear weight, open wound or ulcer, fever, redness spreading up the leg, or sudden severe pain with swelling
  \item Accompanied by other concerning signs
\end{itemize}

\section*{Related Symptoms}
\begin{itemize}
  \item Heel pain
  \item Ankle pain
  \item Toe pain
  \item Swelling
  \item Numbness or tingling
\end{itemize}

\textbf{Note:} Always consider serious underlying conditions when evaluating persistent foot pain.

\section*{Self-Care / Home Management}
\begin{itemize}
  \item Rest and avoid activities that may aggravate the condition.
  \item Maintain adequate hydration and a balanced diet.
  \item Over-the-counter remedies may provide symptomatic relief (consult a pharmacist or doctor).
  \item Monitor symptoms and seek professional advice if they persist or worsen.
  \item Avoid self-medicating with prescription medications without medical guidance.
\end{itemize}

\section*{Diagnostic Approach}
\begin{itemize}
  \item A thorough history and physical examination are the first steps in evaluation.
  \item Laboratory tests (blood work, urinalysis) may be ordered based on clinical suspicion.
  \item Imaging studies (X-ray, ultrasound, CT, MRI) may be indicated when structural pathology is suspected.
  \item Specialist referral is arranged when the diagnosis remains uncertain or requires specific expertise.
\end{itemize}

\section*{Treatment / Management Overview}
\begin{itemize}
  \item Treatment targets the underlying cause identified through diagnostic evaluation.
  \item Supportive care and symptom management are provided while awaiting definitive therapy.
  \item Medications, lifestyle modifications, and physical therapies may be prescribed as appropriate.
  \item Follow-up ensures treatment effectiveness and allows timely adjustments to the care plan.
\end{itemize}

\clearpage